\hypertarget{realizing-the-rag-system-with-amazon-bedrock-and-aurora-pgvector}{%
\section{Realizing the RAG System with Amazon Bedrock and Aurora
pgvector}}

\textbf{Retrieval-Augmented Generation (RAG)} is a method that combines
the power of Large Language Models (LLMs) with a business's proprietary
data. For RAG to work effectively, the heart of the system is the
\textbf{Embedding Model} (to convert text into vectors) and the
\textbf{Vector Database} (to store and search vectors).

Here is how we built this heart entirely using the AWS ecosystem.

\hypertarget{amazon-bedrock-the-power-of-embedding}{%
\subsubsection{1. Amazon Bedrock: The Power of
Embedding}}

Instead of self-hosting models like HuggingFace SentenceTransformer on
EC2 (which is very expensive for GPUs) or Lambda (which runs slowly and
suffers from cold starts), we used \textbf{Amazon Bedrock} with the
\texttt{amazon.titan-embed-text-v2:0} model.

\begin{itemize}
\tightlist
\item
  \textbf{Performance and Cost}: The Bedrock API responds incredibly
  fast, processing embeddings for text chunks in milliseconds. The cost
  is calculated per token, which is extremely cheap for small to medium
  projects.
\item
  \textbf{Absolute Consistency}: A vital rule of RAG is that the data in
  the Database and the user's question (Query) must be converted into
  vectors by the \textbf{same model}. By calling the Bedrock API during
  both the ETL phase and the RAG API phase, we ensure the absolute
  consistency of our vector space.
\end{itemize}

\hypertarget{amazon-aurora-serverless-v2-pgvector}{%
\subsubsection{2. Amazon Aurora Serverless v2 +
pgvector}}

Instead of using third-party Vector Databases like Qdrant, Pinecone, or
Milvus, we decided to use \textbf{Amazon Aurora PostgreSQL} combined
with the \textbf{\texttt{pgvector}} extension.

\begin{itemize}
\tightlist
\item
  \textbf{All-in-one}: PostgreSQL allows storing metadata (like author,
  publish date, URL) in a relational structure (Star Schema) combined
  with a \texttt{vector} column to store embeddings. This makes it easy
  for us to perform hybrid queries: ``Find articles about the economy
  (vector search) but only fetch those published this week (SQL
  filter)''.
\item
  \textbf{HNSW Indexing}: To search quickly across hundreds of thousands
  of data chunks, we configured the \textbf{HNSW (Hierarchical Navigable
  Small World)} index on the vector column. Consequently, cosine
  similarity search queries return results almost instantaneously.
\item
  \textbf{Serverless v2}: The Aurora Database automatically scales up
  and down (from 0.5 to 2 ACUs) based on system load, optimizing costs
  while ensuring high performance when users query the Q\&A system.
\end{itemize}

\hypertarget{conclusion}{%
\subsubsection{Conclusion}}

The architectural combination of Amazon Bedrock and Aurora pgvector
provides a closed-loop (End-to-End) RAG solution entirely on AWS. The
news data never has to leave the system's VPC to make external calls,
ensuring maximum security and reducing network latency. This is the
technical highlight we are most proud of in the entire application
development process.
